\begin{table}[htbp]
\centering
\small
\caption{Exact Shift--Share Decompositions of the Modern PCSI Decline}
\label{tab:decomp}
\begin{threeparttable}
\begin{tabular}{lrr}
\toprule
 & Exclusive primary topic & Deontic $\times$ sanction \\
\midrule
Early period: mean score (1980--1999) & 0.455454 & 0.455454 \\
Late period: mean score (2010--2025) & 0.426638 & 0.426638 \\
\textbf{Observed change} & \textbf{$-0.028816$} & \textbf{$-0.028816$} \\
Within-category contribution & $-0.022229$ & $-0.027077$ \\
Composition contribution & $-0.006587$ & $-0.001739$ \\
Within-category share of decline & 77.1\% & 94.0\% \\
Composition share of decline & 22.9\% & 6.0\% \\
Universities with policies in both periods & 63 & 63 \\
Early and late sentences & 8,353; 28,175 & 8,353; 28,175 \\
\bottomrule
\end{tabular}
\begin{tablenotes}[flushleft]\footnotesize
\item \textit{Notes:} The unit of comparison is a university appearing in both periods; universities are weighted equally. Each sentence is assigned to one primary topic by fixed keyword priority, or to one of eight deontic-by-sanction cells. Category changes are decomposed within each university using midpoint shares and midpoint conditional means, then averaged. If a category is absent in a period, its conditional mean is imputed from the observed period, assigning category entry/exit to composition. The university-level accounting identity holds to machine precision. These alternative decompositions are not additive and are sensitive to the coding rules and missing-category convention. The numbers are descriptive, not causal estimates.
\end{tablenotes}
\end{threeparttable}
\end{table}
